% source: Robin Hartshorne, "Algebraic Geometry", GTM 52 (1977)
% pdf_page: 0316
% doc_page: 299 (Chapter IV, Section 2, "Hurwitz's Theorem")
% retrieved: 2026-07-17
% transcribed_by: claude-fable-5 (vision, rendered at 200dpi via pdftoppm)

% running head: 2 Hurwitz's Theorem / p. 299

\DeclareMathOperator{\Div}{Div}
\DeclareMathOperator{\Pic}{Pic}
% \mathscr (script L, as printed for invertible sheaves) requires mathrsfs.

% Last exercise of the Section 1 exercises (continued from the previous page's
% exercise list).
\begin{exercise}[1.10]
Let $X$ be an integral projective scheme of dimension 1 over $k$, which is
locally complete intersection, and has $p_a = 1$. Fix a point
$P_0 \in X_{\mathrm{reg}}$. Imitate (1.3.7) to show that the map
$P \to \mathscr{L}(P - P_0)$ gives a one-to-one correspondence between the
points of $X_{\mathrm{reg}}$ and the elements of the group $\Pic X$. This
generalizes (II, 6.11.4) and (II, Ex.\ 6.7).
\end{exercise}

\section*{2 Hurwitz's Theorem}

In this section we consider a finite morphism of curves $f\colon X \to Y$,
and study the relation between their canonical divisors. The resulting
formula involving the genus of $X$, the genus of $Y$, and the number of
ramification points is called Hurwitz's theorem.

Recall that the \emph{degree} of a finite morphism $f\colon X \to Y$ of
curves is defined to be the degree $[K(X)\colon K(Y)]$ of the extension of
function fields (II, \S 6).

For any point $P \in X$ we define the \emph{ramification index} $e_P$ as
follows. Let $Q = f(P)$, let $t \in \mathcal{O}_Q$ be a local parameter at
$Q$, consider $t$ as an element of $\mathcal{O}_P$ via the natural map
$f^\#\colon \mathcal{O}_Q \to \mathcal{O}_P$, and define
\[
  e_P = v_P(t),
\]
where $v_P$ is the valuation associated to the valuation ring
$\mathcal{O}_P$. If $e_P > 1$ we say $f$ is \emph{ramified} at $P$, and that
$Q$ is a \emph{branch point} of $f$ (Fig.\ 13). If $e_P = 1$, we say $f$ is
\emph{unramified} at $P$. This definition is consistent with the earlier
definition of unramified (III, Ex.\ 10.3) since our groundfield $k$ is
algebraically closed, and so $k(P) = k(Q)$ for any point $P$ of $X$. In
particular, if $f$ is unramified everywhere, it is \'etale, because in any
case it is flat by (III, 9.7).

% Figure: a wavy curve (labeled X) drawn vertically at left, with a point P
% marked at a local extremum where a dashed horizontal line is tangent to the
% curve; a short horizontal arrow labeled f points right to a vertical line
% (labeled Y) carrying a point Q at the level of the dashed line.
% Caption:
% Figure 13. A finite morphism of curves.

If $\operatorname{char} k = 0$, or if $\operatorname{char} k = p$, and $p$
does not divide $e_P$, we say that the ramification is \emph{tame}. If $p$
does divide $e_P$, it is \emph{wild}.

Recall that we have defined a homomorphism
$f^*\colon \Div Y \to \Div X$ of the groups of divisors, by setting
\[
  f^*(Q) = \sum_{P \to Q} e_P \cdot P
\]
for any point $Q$ of $Y$, and extending by linearity (II, \S 6). If $D$ is a
divisor on $Y$, then $f^*(\mathscr{L}(D)) \cong \mathscr{L}(f^*D)$
(II, Ex.\ 6.8), so that this $f^*$ on divisors is compatible with the
homomorphism $f^*\colon \Pic Y \to \Pic X$ on invertible sheaves.
